\section{Conclusion}
\label{sec:conclusion}
Research on generative graphic design encompasses flat image synthesis alongside spatial layouts, vector paths, compositional layers, and executable design documents, rather than following a single progression between output types.
These advances demand more discriminating evidence than a binary editability label.
We introduced a representation--editability--evaluation crosswalk and used representative research to separate delivered structure, supported manipulation, and measured task outcomes.
This view makes it possible to recognize the achievements of layout, raster, vector, typography, and agentic systems without conflating their metrics.
The synthesis identifies a recurring evidence gap between an artifact's editable structure, a demonstrated operation, and quantitative performance on that operation.
A representation-aware benchmark should therefore control the information supplied to each method, distinguish permitted reflow from protected content, and report target correctness, artifact validity, visual quality, and effort separately.
The literature audit remains provisional; this manuscript does not claim a completed systematic screening process or a unified empirical ranking across incompatible tasks.

\paragraph{\ac{llm} usage statement.}
\ac{ai} assistance was used to prepare this working manuscript and initial bibliography.
The authors remain responsible for checking citation accuracy, interpretations, and source-level experimental claims before formal submission.
